\begin{abstract}
LLM agent skills (folders of instructions, scripts and reference files that an agent loads on
demand) are becoming commercial products, yet they ship as plain text that anyone can copy. We
present \sys{}, a packaging scheme in which a skill is encrypted with HPKE to a public key that an
AI provider publishes at a well-known URL on its own domain. The owner seals offline with no
handshake, packages can be distributed through untrusted channels, and only the provider's servers
can open them; clients send a reference and a licence and never receive the skill. Our prototype
seals and opens a 10\,KB skill in under half a millisecond. We then measure the channel encryption
cannot close: the model must read the skill, and its answers go to the user. Across eight
synthetic skills, 24 extraction attacks and two open 7--8B models, 72--75\% of attack
conversations disclose at least half of a skill's rules and 35--41\% reproduce it nearly verbatim.
A confidentiality instruction reduces substantial extraction to 19\% and 46\%, but for one model
at the cost of failing a third of legitimate tasks, mostly by refusing them. A de-obfuscating lexical output filter
reduces substantial extraction to 1.6--3.1\% but blocks 12.5--25\% of benign replies, because
models quote rules while applying them; embedding-similarity filtering is far weaker. Per-customer
canaries appear in 98\% of verbatim leaks, enabling attribution. In a tool-use case study, scripts
co-located with a code interpreter are extractable, while an owner-hosted tool protects the code
but still reveals its rate table to ordinary queries. Sealing protects a skill's artifacts from
copying; it does not make its content or function secret.
\end{abstract}
